% Generated draft: requires mathematical and source review.
\documentclass[11pt]{article}
\usepackage[margin=25mm]{geometry}
\usepackage{fontspec}
\setmainfont{Latin Modern Roman}
\usepackage{amsmath,amssymb,mathtools,unicode-math}
\setmathfont{Latin Modern Math}
\usepackage{enumitem,needspace,xurl,hyperref}
\setlength{\parindent}{0pt}
\setlength{\parskip}{5pt}
\setlength{\emergencystretch}{3em}
\newcommand{\N}{\mathbb{N}}
\newcommand{\Z}{\mathbb{Z}}
\newcommand{\Q}{\mathbb{Q}}
\newcommand{\R}{\mathbb{R}}
\newcommand{\C}{\mathbb{C}}
\newcommand{\F}{\mathbb{F}}
\newcommand{\A}{\mathbb{A}}
\newcommand{\PP}{\mathbb{P}}
\newcommand{\E}{\mathbb{E}}
\newcommand{\eps}{\varepsilon}
\newcommand{\dd}{\,\mathrm{d}}
\begin{document}
\section*{UnsolvedMath ID: 30003342}
\textbf{Source problem code:} catalogue entry
\subsection{Definitions and mathematical statement}
The supplied record concerns women mathematicians in twentieth-century Denmark. It gives a historical, biographical, and sociological research question, rather than a mathematical conjecture. No mathematical objects, hypotheses, or proposition suitable for formalization are specified. Accordingly, the research question is retained without introducing mathematical quantifiers or an artificial mathematical model.

For women mathematicians in twentieth-century Denmark, investigate the following questions:
\begin{enumerate}
\item How did they enter university positions?
\item What motivated their mathematical interests and their choice of research area?
\item What was their daily departmental life like?
\item How did they balance work and family?
\item What reactions did they experience as some of the few women?
\end{enumerate}

The source does not specify the population boundaries, the relevant institutions, the meaning of being one of the few women, or a criterion for a complete comparative answer. These qualifications remain unresolved. The supplied statement-verification note classifies the item as a historical, biographical, or sociological research prompt rather than an open mathematical problem; its statement status is recorded as \texttt{unrecoverable}, and no separate clean mathematical statement is supplied. The registry's \texttt{open} label is retained as supplied curation, not as an independently verified mathematical status.
\subsection{Short English statement}
This historical research prompt asks about university careers, mathematical motivations, departmental life, work and family, and reactions experienced by women mathematicians in twentieth-century Denmark. It supplies no mathematical proposition to formalize; its population boundaries and standards for a complete answer require clarification.
\subsection{Sources}
\begin{itemize}
\item Local UnsolvedMath registry, permanent ID \(30003342\), title: Careers of Danish Women Mathematicians; source code: \texttt{catalogue entry}; collection: Oberwolfach Reports Open Problems. The supplied proposed year is 2017; no individual proposer is identified. The original record is retained for provenance, with no separate clean statement promoted.
\item Supplied source citation: Mini-Workshop: Women in Mathematics: Historical and Modern Perspectives (2018), Oberwolfach Reports, pp.\ 83--131, DOI 10.4171/owr/2017/2. \url{https://doi.org/10.4171/owr/2017/2}. Also listed among the supplied literature sources as ``Oberwolfach report on women in mathematics.'' The citation's 2018 designation differs from the supplied proposed year and year tag, both 2017; this discrepancy is unresolved.
\item Supplied literature source: Against All Odds: Women's Ways to Mathematical Research. \url{https://doi.org/10.1007/978-3-030-47610-6}. No authors or specific supporting passages are supplied.
\item Supplied literature assessment dated exactly 2026-08-21: status \texttt{open}. The assessment states that the cited interviews and later historical work recover parts of individual careers, motivations, and work-family experiences, but do not provide a complete evidence-based answer to the multipart comparative question. This is supplied dated curation derived from the OWR record, not independently verified literature assessment. Time-sensitive status and cited claims require independent verification.
\end{itemize}
\end{document}
